%Version 3 October 2023
% See section 11 of the User Manual for version history
%
%%%%%%%%%%%%%%%%%%%%%%%%%%%%%%%%%%%%%%%%%%%%%%%%%%%%%%%%%%%%%%%%%%%%%%
%%                                                                 %%
%% Please do not use \input{...} to include other tex files.       %%
%% Submit your LaTeX manuscript as one .tex document.              %%
%%                                                                 %%
%% All additional figures and files should be attached             %%
%% separately and not embedded in the \TeX\ document itself.       %%
%%                                                                 %%
%%%%%%%%%%%%%%%%%%%%%%%%%%%%%%%%%%%%%%%%%%%%%%%%%%%%%%%%%%%%%%%%%%%%%

%%\documentclass[referee,sn-basic]{sn-jnl}% referee option is meant for double line spacing

%%=======================================================%%
%% to print line numbers in the margin use lineno option %%
%%=======================================================%%

%%\documentclass[lineno,sn-basic]{sn-jnl}% Basic Springer Nature Reference Style/Chemistry Reference Style

%%======================================================%%
%% to compile with pdflatex/xelatex use pdflatex option %%
%%======================================================%%

%%\documentclass[pdflatex,sn-basic]{sn-jnl}% Basic Springer Nature Reference Style/Chemistry Reference Style


%%Note: the following reference styles support Namedate and Numbered referencing. By default the style follows the most common style. To switch between the options you can add or remove “Numbered” in the optional parenthesis. 
%%The option is available for: sn-basic.bst, sn-vancouver.bst, sn-chicago.bst%  
 
\documentclass[sn-nature]{sn-jnl}% Style for submissions to Nature Portfolio journals
%%\documentclass[sn-basic]{sn-jnl}% Basic Springer Nature Reference Style/Chemistry Reference Style
%%\documentclass[sn-mathphys-num]{sn-jnl}% Math and Physical Sciences Numbered Reference Style 
%%\documentclass[sn-mathphys-ay]{sn-jnl}% Math and Physical Sciences Author Year Reference Style
%%\documentclass[sn-aps]{sn-jnl}% American Physical Society (APS) Reference Style
%%\documentclass[sn-vancouver,Numbered]{sn-jnl}% Vancouver Reference Style
%%\documentclass[sn-apa]{sn-jnl}% APA Reference Style 
%%\documentclass[sn-chicago]{sn-jnl}% Chicago-based Humanities Reference Style

%%%% Standard Packages
%%<additional latex packages if required can be included here>

\usepackage{graphicx}%
\usepackage{multirow}%
\usepackage{amsmath,amssymb,amsfonts}%
\usepackage{amsthm}%
\usepackage{mathrsfs}%
\usepackage[title]{appendix}%
\usepackage{xcolor}%
\usepackage{textcomp}%
\usepackage{manyfoot}%
\usepackage{booktabs}%
\usepackage{algorithm}%
\usepackage{algorithmicx}%
\usepackage{algpseudocode}%
\usepackage{listings}%
\usepackage{subcaption}
\usepackage{comment}
%%%%

%%%%%=============================================================================%%%%
%%%%  Remarks: This template is provided to aid authors with the preparation
%%%%  of original research articles intended for submission to journals published 
%%%%  by Springer Nature. The guidance has been prepared in partnership with 
%%%%  production teams to conform to Springer Nature technical requirements. 
%%%%  Editorial and presentation requirements differ among journal portfolios and 
%%%%  research disciplines. You may find sections in this template are irrelevant 
%%%%  to your work and are empowered to omit any such section if allowed by the 
%%%%  journal you intend to submit to. The submission guidelines and policies 
%%%%  of the journal take precedence. A detailed User Manual is available in the 
%%%%  template package for technical guidance.
%%%%%=============================================================================%%%%

%% as per the requirement new theorem styles can be included as shown below
\theoremstyle{thmstyleone}%
\newtheorem{theorem}{Theorem}%  meant for continuous numbers
%%\newtheorem{theorem}{Theorem}[section]% meant for sectionwise numbers
%% optional argument [theorem] produces theorem numbering sequence instead of independent numbers for Proposition
\newtheorem{proposition}[theorem]{Proposition}% 
%%\newtheorem{proposition}{Proposition}% to get separate numbers for theorem and proposition etc.

\theoremstyle{thmstyletwo}%
\newtheorem{example}{Example}%
\newtheorem{remark}{Remark}%

\theoremstyle{thmstylethree}%
\newtheorem{definition}{Definition}%

%% New commands
\newcommand{\prob}{\text{Pr}}

\raggedbottom
%%\unnumbered% uncomment this for unnumbered level heads

\begin{document}

\title[Article Title]{Identification of mechanistic heterogeneity across thousands of rare and complex traits}

%%=============================================================%%
%% GivenName	-> \fnm{Joergen W.}
%% Particle	-> \spfx{van der} -> surname prefix
%% FamilyName	-> \sur{Ploeg}
%% Suffix	-> \sfx{IV}
%% \author*[1,2]{\fnm{Joergen W.} \spfx{van der} \sur{Ploeg} 
%%  \sfx{IV}}\email{iauthor@gmail.com}
%%=============================================================%%

\author*[1,2]{\fnm{First} \sur{Author}}\email{iauthor@gmail.com}

\author[2,3]{\fnm{Second} \sur{Author}}\email{iiauthor@gmail.com}
\equalcont{These authors contributed equally to this work.}

\author[1,2]{\fnm{Third} \sur{Author}}\email{iiiauthor@gmail.com}
\equalcont{These authors contributed equally to this work.}

\affil*[1]{\orgdiv{Department}, \orgname{Organization}, \orgaddress{\street{Street}, \city{City}, \postcode{100190}, \state{State}, \country{Country}}}

\affil[2]{\orgdiv{Department}, \orgname{Organization}, \orgaddress{\street{Street}, \city{City}, \postcode{10587}, \state{State}, \country{Country}}}

\affil[3]{\orgdiv{Department}, \orgname{Organization}, \orgaddress{\street{Street}, \city{City}, \postcode{610101}, \state{State}, \country{Country}}}

%%==================================%%
%% Sample for unstructured abstract %%
%%==================================%%

\abstract{TODO
}

%%================================%%
%% Sample for structured abstract %%
%%================================%%

% \abstract{\textbf{Purpose:} The abstract serves both as a general introduction to the topic and as a brief, non-technical summary of the main results and their implications. The abstract must not include subheadings (unless expressly permitted in the journal's Instructions to Authors), equations or citations. As a guide the abstract should not exceed 200 words. Most journals do not set a hard limit however authors are advised to check the author instructions for the journal they are submitting to.
% 
% \textbf{Methods:} The abstract serves both as a general introduction to the topic and as a brief, non-technical summary of the main results and their implications. The abstract must not include subheadings (unless expressly permitted in the journal's Instructions to Authors), equations or citations. As a guide the abstract should not exceed 200 words. Most journals do not set a hard limit however authors are advised to check the author instructions for the journal they are submitting to.
% 
% \textbf{Results:} The abstract serves both as a general introduction to the topic and as a brief, non-technical summary of the main results and their implications. The abstract must not include subheadings (unless expressly permitted in the journal's Instructions to Authors), equations or citations. As a guide the abstract should not exceed 200 words. Most journals do not set a hard limit however authors are advised to check the author instructions for the journal they are submitting to.
% 
% \textbf{Conclusion:} The abstract serves both as a general introduction to the topic and as a brief, non-technical summary of the main results and their implications. The abstract must not include subheadings (unless expressly permitted in the journal's Instructions to Authors), equations or citations. As a guide the abstract should not exceed 200 words. Most journals do not set a hard limit however authors are advised to check the author instructions for the journal they are submitting to.}

\keywords{keyword1, Keyword2, Keyword3, Keyword4}

%%\pacs[JEL Classification]{D8, H51}

%%\pacs[MSC Classification]{35A01, 65L10, 65L12, 65L20, 65L70}

\maketitle

\section{Introduction}\label{sec:introduction}

Identifying genes, pathways, and mechanisms that underlie human disease is central to understanding disease etiology[], developing new treatments[], and advancing precision medicine[]. Genome wide association studies (GWAS) have produced common variant associations for tens of thousands of complex traits[], each of which in principle is translatable to effects on genes and downstream pathways[]. In practice, this translation is difficult because most associations are due to variants with unknown effects on a complicated regulatory network[] which, for complex multifactorial disorders like type 2 diabetes[] or coronary artery disease[], impacts a heterogeneous collection of molecular and cellular mechanisms. Detailed investigation has made it possible to, for some traits, organize a subset of associations into novel disease pathogeneses[], gene programs[], or disease subtype-specific risk factors[]. Considering GWAS as a research field, however, it is still debated whether most associations will cohere within biological pathways[] or instead implicate nearly every gene expressed in disease-relevant cell types[]. 

While experimental work is usually needed to uncover mechanisms underlying GWAS associations, computational methods can help predict which genes and pathways merit further study. Most research to date has focused on predicting which gene, within a genome-wide significant associated locus, mediates the effects of the causal variant[]. These “effector gene” prediction algorithms are broadly classifiable into “SNP-to-gene” approaches[] that use genetic and genomic features to link causal variants to genes and “convergence-based” approaches[] that use annotations of gene functions (also representable as gene sets) to prioritize among genes nearby an association. While knowledge of an effector gene is a valuable first step toward elucidating a disease mechanism, SNP-to-gene approaches provide no information about affected downstream pathways, and convergence-based approaches cannot yet synthesize the gene annotations they use into mechanistic predictions due to the lack of interpretability of their model parameters[]. Large language models (LLMs) may one day complement – or even replace[] – these approaches with mechanistic predictions, but their models are text-based and cannot statistically analyze genetic or genomic data specific to a GWAS.

Furthermore, substantial information about disease mechanisms can be extracted from associations that do not reach genome-wide significance (i.e. for which effector gene predictions are inapplicable). In related settings, methods that use polygenic signal have more power than methods that consider only significant associations to identify disease-relevant epigenomic annotations[], cell types[], or pathways[]. Moreover, many if not the majority of experimentally studied candidate disease genes do not originate from GWAS associations[]. Being able to provide estimates of “genetic support” for any disease and any candidate gene in the genome – as well as potential disease-relevant pathways or mechanisms in which the gene may participate – could significantly increase the success rate of therapeutic targets[] and guide experimental resources to where they are most valuable[].

In this study, we introduce a new approach – termed PIGEAN – to identify mechanisms underlying complex traits by jointly estimating gene and gene set associations from GWAS summary statistics. We show how this approach probabilistically generalizes both leading methods for convergence-based effector gene prioritization and leading methods for gene set enrichment analysis, enabling them to be combined to produce interpretable and comparable values for every gene in the genome – most notably the vast majority of genes that lie outside of genome-wide significant loci. The result is a quantitative estimate of genetic support for each gene, interpretable as the probability that the gene participates in a pathway involved in disease. This estimate is decomposable into “direct” genetic support – which quantifies the probability that disease-associated variants exert their effects on the gene – and “indirect” genetic support – which quantifies the probability that disease-associated variants exert their effects on other genes within the same pathway as the gene. Correspondingly, each gene set receives an estimate regarding the effect that it has on the (log-odds) indirect genetic support of each gene in the set. With access to these jointly estimated gene and gene set associations for a disease, we show that it is possible to identify them as proxies for higher-order latent disease mechanisms (defined as weighted collections of genes and gene sets). Applying this approach to 1,112 traits, we show that associations for nearly all well-powered GWAS are decomposable into distinct, coherent, and interpretable mechanisms, reaffirming the view that –  contrary to recurring fears[] – GWAS are a reliable route to discover valuable insights into core disease biology.

\section{Methods}\label{sec:methods}
TODO.

\section{Results}\label{sec:results}
TODO.

%%% NOTES
\begin{comment}

We use the jointly inferred genes and gene sets to decompose GWAS into distinct mechanisms, each characterized by weighted genes and gene sets. Traits with extreme mechanistic heterogeneity, such as WHR, SBP, Aging, and CAD decompose cleanly. These results indicate that, despite some fears to the contrary, GWAS clearly localize to pathways provided with appropriate statistical analysis.

Architectural comparison across traits

Correlate trait x trait matrix and compare genes to genetic correlations. Also use gene set values here
PCs, how many independent traits?
Find diseases *more* similar in pathway space than gene space, or *more* similar in gene space than variant space.
- There should be a variant here (based on genes), but there isn't
- There should be a gene here (based on pathways), but there isn't

Do types of features that matter vary by trait? Does it depend on measurable properties of traits such as organ system or tissues from LDSC?

Do the inferences change if we add exome sequence data?
Are exome pathways similar to GWAS pathways?

Which common disease maps best to each rare disease? Use Orphanet gene sets

Different genetic architecture of different traits. Do some traits have more direct than indirect?

Which disease has the most similar mouse phenotypes? That is, some traits will have top gene sets be mouse while others will have MSigDB. Any pattern here?
Which diseases have more mouse vs. more MSigDB vs. more Text etc

Going off of previous methods (signals + Pops or depict) how many signals are explained? What about now? May be a useful statistic in abstract (we show that, whole previous methods estimate only 20%, actually closer to 80%)

Paragraph in paper on noise with fixed p value threshold. Well known but here very explicit. Less clean results. Confounders play a role. Confounders related to effect sizes, does that imply the omnigenic trans effects?

Potential pathways and mechanisms across 1000 traits

Novel gene sets significant after iterating
Gene sets that are not significant after iterating
Novel gene sets relative to MAGMA

These gene sets seem like they are relevant. But actually these other gene sets are better .

How many times is the top gene set a good mouse gene set?

Which gene set is most pleitropic?


\begin{enumerate}
    \item Make a plot of gwas sample size vs. sum of gene set betas
    \item Make a plot of sum of direct gene probabilities vs sum of gene set betas
\end{enumerate}


Fine mapping gene sets. Compare with gene sets prioritized by other methods. Do some gene sets explain others away (KEGG, text terms from PubMed)? 

Novel mechanisms, including clustering and prediction 

Can we show that the later GWAS loci are crappier? Is it real that detect power works so much better? Divide GWAS into SNPs with p<5e-30 and those with p > 5e-30

Retinoic acid and oily fish
Neuroticism and copy number syndrome
Fat intake and lipid genes
Myopathy genes and heart structure traits (top 5  and more of Alan's genes all have cardiovascular as top disease group)

Segmentation of GWAS into distinct mechanisms

Apply clustering. Show associations with distinct trait patterns

Longevity parsing mechanisms

Renal cluster for t2d -- including HNF1B

Partitioned polygenic scores based on pathways
Pathway partitioning of signals, do you see different (a) epigenomic or (b) trait signatures?
Look at Kellis SCZ dissection paper and see if we can match these to genes somehow
Factor the GWAS. Show this improves PRS and tissues. Then see if have novel mechanisms. Use GPT to label the gene sets. Has it been reported before?

Partitioned rare variant risk scores based on the pathways – HbA1C and T2D
Overall rare variant risk scores based on the top ranked genes

Application: for those traits we segment, how accurate were we compared to custom work done for the GWAS?

SBP is very interesting
CAD is very interesting
WHR is very interesting
BMI okay
Multi-variate aging mechanisms

NPY – looks like mechanisms vary based on which gene sets you throw into the model. Some cognitive, some metabolic. Recent Nature paper on NPY may support this

Different mechanisms from FIadjBMI and FI, WHR and WHRadjBMI, other adjustments

LDL less heterogeneity than Height or SBP: going to much larger gene set lists helps here, whereas for the traits like SBP maybe too much
Which traits have least clusters? Which have most? Does this relate to the genetic architecture? Or something phenotypic? 
Do less heterogeneous traits have stronger rare variant associations?
Does it track with number of genes with combined above the threshold?

How many / what fraction of genes with known mechanisms are within the clusters? Does this vary across traits?

Compare to CAD endothelial gene programs
Just out of the box gene sets
With full matrix from Jesse’s group

Criticism may be that we can only detect known biology.
Look at genes that are not placed into a factor (presumably high direct but low indirect support)
See how much is known about them (how many annotations are there).
Are there a lot of annotations, but no disease-relevant ones -> may support omnigenic model
Are there few annotations -> should focus on studying these genes
Are the genes without annotations likely to be discovered more recently -> newer GWAS won’t identify new mechanisms
Do we discover new mechanisms over time?
Do things get more coherent over time?

Take Orphanet gene sets. Do basic enrichment for each GWAS, just combined scores. Then, do enrichment for each mechanism (either combined or combined * lambda). Show greater enrichment for monogenic disease genes within more homogeneous mechanisms

Partitioned polygenic scores: we should generate PRS for each cluster, and see if they stratify patients


Rare disease pheWAS


Factoring trait phewas x all



Compare mechanisms of genes in GWAS loci to mechanisms outside

Architecture of mechanisms and coherence of GWAS
Plot number of mechanisms or number of gene sets over time
Show increases
GWAS are converging to pathways as sample sizes increase 
Traits with many associations and mechanisms, many and no mechanisms
Many mechanisms could mean number but also sum of gene scores or max gene score
Scatter plot of num associations vs mechanisms. Which traits are top and bottom? Most coherent vs least?
Look at traits with a lot fewer gene set associations and mechanisms that inspected based on the number of associations 
Then to the reversed: a lot of mechanisms with only a small number of associations 
Intelligence is one example

CAN WE:
Find additional genes that match a mechanism (even those without combined score above 1)?
Partition heritability across mechanisms
Ask about mapping a mechanism into the set or genes/gene sets? E.g. take the handwritten mechanisms from the recent T1D/T2D Cell review and somehow find them?
Either match them to the mechanisms we’ve inferred, or they upload a set of genes or gene sets and see where they match?


\end{comment}


%%% Supplementary Materials

\begin{appendices}

\section{Supplementary Information}\label{sec:appendix-supplementary-info}

This appendix contains supplementary information that is not an essential part of the text itself but which may be helpful in providing a more comprehensive understanding of the research problem or it is information that is too cumbersome to be included in the body of the paper.

%%=============================================%%
%% For submissions to Nature Portfolio Journals %%
%% please use the heading ``Extended Data''.   %%
%%=============================================%%

%%=============================================================%%
%% Sample for another appendix section			       %%
%%=============================================================%%

%% \section{Example of another appendix section}\label{secA2}%
%% Appendices may be used for helpful, supporting or essential material that would otherwise 
%% clutter, break up or be distracting to the text. Appendices can consist of sections, figures, 
%% tables and equations etc.

\end{appendices}

%%===========================================================================================%%
%% If you are submitting to one of the Nature Portfolio journals, using the eJP submission   %%
%% system, please include the references within the manuscript file itself. You may do this  %%
%% by copying the reference list from your .bbl file, paste it into the main manuscript .tex %%
%% file, and delete the associated \verb+\bibliography+ commands.                            %%
%%===========================================================================================%%

\bibliography{references}% common bib file
%% if required, the content of .bbl file can be included here once bbl is generated
%%\input sn-article.bbl


\end{document}
